% !TeX root = main.tex
\chapter{群}

% 半群的定义
% 幺半群的定义
% 幺元的性质：要么左幺等于右幺，且幺元唯一，要么存在（多个）左幺元或右幺元

% 群的定义
% 逆元的性质：要么左逆等于右逆，且逆元唯一，要么存在（多个）左逆元或右逆元
% 结合律的作用：规定了唯一形式的一种乘积，因此可定义幂运算（an = a · a · ... · a）a^-n 是良定义的 a^0 = e（0a）是一种规定
% 交换群etc 用+表示运算

% 例：Zn + 的定义
% Zn * 的定义（严格说明）
% 贝祖定理及拓展欧几里得算法
% wallis 定理：如果p是素数，那么对于任意整数a，a^p ≡ a (mod p) 群伦证明

% 子群
% 子群的逆和幺元都和大群里的一样
% 定理：关于运算封闭等价 ab-1 在H中

% orders of elements
% 定义
% ab的阶、ord(a^k) = ord(a) / gcd(ord(a), k)

% 群同态/同构

\section{循环群}
\begin{theorem}[循环群的结构]
    设\(G\) 是循环群：
    \begin{itemize}
        \item 若\(\abs{G}\)是有限的，则\(G\)同构于\(\mathbb{Z}_n\)，其中\(n = \abs{G}\)；
        \item 若\(\abs{G}\)是无限的，则\(G\)同构于\(\mathbb{Z}\)。
    \end{itemize}
\end{theorem}
\begin{theorem}[循环群子群结构]
    循环群的任意子群也是循环群。
    若\(m|n\), 则\(\mathbb{Z}_n\)有且仅有一个子群同构于\(\mathbb{Z}_m\)。
\end{theorem}